% !TeX root = ../main.tex
%%% ---------------------------------------------------------------------------
%%% 实验原理
%%%
%%% 写作提示：原理部分要写到**能据此复现**，但不要抄教材。
%%% 判断标准是：读者看完这一节，能说出你为什么这么设计实验。
%%% 公式的正斜体规则见 ../../STANDARDS.md §1.3：
%%%   变量斜体、常数 \e 正体、微分 \dd 正体、说明性下标 \subtext{} 正体。
%%% ---------------------------------------------------------------------------
\section{实验原理}
\label{sec:principle}

\subsection{优化器的更新规则}
\label{sec:principle:update}

设第 $t$ 步的参数为 $\params_t$，损失函数为 $\loss(\params)$，梯度
$\vect{g}_t = \nabla_{\params} \loss(\params_t)$。

带动量的随机梯度下降按

\begin{equation}
  \label{eq:sgdm}
  \vect{v}_t = \mu \vect{v}_{t-1} + \vect{g}_t, \qquad
  \params_{t+1} = \params_t - \lr \vect{v}_t
\end{equation}

更新，其中 $\mu$ 为动量系数。Adam 引入一阶与二阶矩的指数滑动平均：

\begin{equation}
  \label{eq:adam}
  \params_{t+1} = \params_t
    - \lr \frac{\hat{\vect{m}}_t}{\sqrt{\hat{\vect{v}}_t} + \epsilon},
  \qquad
  \hat{\vect{m}}_t = \frac{\vect{m}_t}{1 - \beta_1^{\,t}},\quad
  \hat{\vect{v}}_t = \frac{\vect{v}_t}{1 - \beta_2^{\,t}}
\end{equation}

AdamW 与 Adam 的区别在于把权重衰减从梯度中剥离，直接作用于参数：

\begin{equation}
  \label{eq:adamw}
  \params_{t+1} = \params_t
    - \lr \left( \frac{\hat{\vect{m}}_t}{\sqrt{\hat{\vect{v}}_t} + \epsilon}
                 + \lambda \params_t \right).
\end{equation}

\cref{eq:adam} 与\cref{eq:adamw} 的差别看似微小，但在正则化较强时会改变
有效学习率的尺度，这正是本实验要测量的对象之一\cite{loshchilov2019adamw}。

\subsection{重复实验与不确定度}
\label{sec:principle:uncertainty}

神经网络训练受随机初始化、数据打乱顺序、GPU 非确定性算子等因素影响，
单次结果不足以支撑结论。对同一配置独立重复 $n$ 次，得到观测值
$\setof{x_i}_{i=1}^{n}$，则

\begin{equation}
  \label{eq:mean-std}
  \bar{x} = \frac{1}{n}\sum_{i=1}^{n} x_i, \qquad
  s = \sqrt{\frac{1}{n-1}\sum_{i=1}^{n}\left(x_i - \bar{x}\right)^2},
\end{equation}

平均值的标准不确定度（A 类评定）为

\begin{equation}
  \label{eq:uncertainty}
  u(\bar{x}) = \frac{s}{\sqrt{n}}.
\end{equation}

报告结果时写作 $\bar{x} \pm u(\bar{x})$，有效数字保留到不确定度的第一位
非零数字（GB/T 3101 与 JJF 1059.1）。

\begin{caution}
  只报单次最好结果（cherry-picking）是实验报告最常见的失分点。
  即使差异明显，也要给出重复次数与波动范围，否则读者无法判断结论是否可靠。
\end{caution}
